\documentclass[11pt,a4paper]{article}
\usepackage{geometry}
\usepackage{indentfirst}

% === DYNAMIC PREAMBLE ===
\usepackage{polyglossia}
\setdefaultlanguage{english}

        \usepackage{fontspec}
        \setmainfont{Times New Roman}[
            Path=../../../Fonts/,
            UprightFont    = * ,
            BoldFont       = *-Bold ,
            ItalicFont     = *-Italic ,
            BoldItalicFont = *-BoldItalic
        ]
% === END DYNAMIC PREAMBLE ===

\usepackage{fancyhdr}
\usepackage{setspace}
\usepackage{tabularx}
\usepackage{booktabs}
\usepackage{framed}
\usepackage{xcolor}
\usepackage{hyperref}

\setlength{\parindent}{3em}

\geometry{
    top=1in,
    bottom=1.5in,
    left=1in,
    right=0.8in,
    footskip=60pt
}

% === LINK COLORS ===
\hypersetup{
    colorlinks=true,
    linkcolor=black,
    urlcolor=black,
    citecolor=black
}

% ==========================================
% FONT COMMANDS
% ==========================================
\newcommand{\headerfont}{\sffamily\fontsize{8}{12}\selectfont}
\newcommand{\contentfont}{\fontsize{11}{13}\selectfont}
\newcommand{\titlefont}{\fontsize{17}{22}\bfseries\selectfont}
\newcommand{\subheadingfont}{\fontsize{11}{13}\bfseries\selectfont}
\newcommand{\footerfont}{\fontsize{8}{12}\selectfont}
\newcommand{\citationfont}{\fontsize{8}{12}\selectfont}
\newcommand{\referencefont}{\fontsize{9}{11}\selectfont}

% ==========================================
% HEADER & FOOTER
% ==========================================
\pagestyle{fancy}
\fancyhf{}

\fancyhead[L]{\headerfont Journal of AI Research}
\fancyhead[R]{\headerfont Volume 06, Issue 02, 2026}
\renewcommand{\headrulewidth}{0pt}


\fancyfoot[C]{%
    \begin{minipage}{\textwidth}
        {\footerfont \makebox[0pt][l]{}\hfill \thepage \hfill \makebox[0pt][r]{2024-001}}\\[0.8em]
        \raggedright
        \citationfont
        \begin{tabularx}{\textwidth}{@{}l X@{}}
            \textbf{Citation Information:} & Smith J. (17). Machine Learning: Transforming Science, Business, and Life. \textit{Journal of AI Research}, 06(02), 56.
        \end{tabularx}
    \end{minipage}%
}
\renewcommand{\footrulewidth}{0pt}

% ==========================================
% MAIN DOCUMENT
% ==========================================
\begin{document}
\contentfont

\begin{center}
    {\titlefont Machine Learning: Transforming Science, Business, and Life\par}

    \vspace{1.2em}

    {\fontsize{13}{10}\selectfont\noindent\textbf{John Smith*, Computer Science}}

    \vspace{1.2em}

    {\subheadingfont ABSTRACT}
\end{center}

\vspace{0.2em}

\begingroup
\itshape
\begin{otherlanguage}{english}
This document synthesizes research on the expansive applications of machine learning across various domains. It highlights ML's role in accelerating scientific discovery, including drug development, materials science, and climate modeling. The text details advancements in natural language processing, cybersecurity, healthcare, and robotics driven by ML. Furthermore, applications in financial forecasting, intelligent transportation systems, agriculture, and renewable energy are discussed, underscoring ML's transformative potential.
\end{otherlanguage}
\par
\endgroup

\vspace{1em}

\noindent\textbf{Keywords:} Machine Learning; Artificial Intelligence; Scientific Discovery; Natural Language Processing; Cybersecurity; Healthcare; Robotics; Financial Forecasting; Intelligent Transportation Systems; Agriculture

\vspace{1em}

% ==========================================
% INTRODUCTION
% ==========================================
\begin{center}
    {\subheadingfont INTRODUCTION}
\end{center}

\vspace{0.2em}

\begin{otherlanguage}{english}
The pervasive integration of machine learning (ML) into scientific endeavors is fundamentally reshaping the landscape of research and discovery across numerous disciplines. This transformative technology is proving instrumental in accelerating the pace of innovation, refining experimental methodologies, and unlocking unprecedented capabilities in data analysis.
\end{otherlanguage}

\vspace{1em}

% ==========================================
% DESCRIPTION
% ==========================================
\begin{center}
    {\subheadingfont DESCRIPTION}
\end{center}

\vspace{0.2em}

\begin{otherlanguage}{english}
The application of machine learning (ML) in scientific research offers a paradigm shift, enabling the exploration of complex phenomena and the acceleration of discovery processes.
\end{otherlanguage}

\vspace{1em}

% ==========================================
% DISCUSSION
% ==========================================
\begin{center}
    {\subheadingfont DISCUSSION}
\end{center}

\vspace{0.2em}

\begin{otherlanguage}{english}
The collective insights presented underscore the profound and pervasive influence of machine learning across a multitude of sectors.
\end{otherlanguage}

\vspace{1em}

% ==========================================
% CONCLUSION
% ==========================================
\begin{center}
    {\subheadingfont CONCLUSION}
\end{center}

\vspace{0.2em}

\begin{otherlanguage}{english}
This collection of research explores the widespread applications of machine learning across diverse fields.
\end{otherlanguage}

\vspace{1em}

% ==========================================
% REFERENCES
% ==========================================
\begin{center}
    {\subheadingfont REFERENCES}
\end{center}

\vspace{0.2em}

{\referencefont
    \noindent\hangindent=3em\hangafter=1
            Marcos VR, Sergej FGW, Alana LVJ et al.
        (2021).
            \href{https://pubmed.ncbi.nlm.nih.gov/33753790/}{Machine learning for scientific discovery.}
        \textit{Nature},
            591(591):523-530.
    
    \noindent\hangindent=3em\hangafter=1
            Quan L, Yongliang S, Wenfeng C et al.
        (2022).
            \href{https://ieeexplore.ieee.org/abstract/document/9530071/}{A Survey of Deep Learning for Natural Language Processing.}
        \textit{IEEE Transactions on Pattern Analysis and Machine Intelligence},
            44(12):8275-8291.
    
}

\vfill
\noindent
\begin{framed}
\fontsize{8}{10}\selectfont
\noindent
\textbf{Received:} 18-Jul-2026, Manuscript No. 2024-001;
\textbf{Editor assigned:} 20-Jul-2026, Pre QC No. 2024-001(PQ);
\textbf{Reviewed:} 03-Aug-2026, QC No. 2024-001;
\textbf{Revised:} 10-Aug-2026, Manuscript No. 2024-001(R);
\textbf{Published:} 2026-08-17
\end{framed}

\end{document}